\chapter{Implementation Architecture}
\label{chap:architecture}

\section{From monolithic contracts to specialized contracts}

The architectural work of the project followed a useful negative result: adding all variants
to one contract was not a good gas optimization strategy. The first monolithic variant implementation
extended the existing optimistic and dispute contracts with additional flags for
\texttt{no\_S\_deposit}, \(S=B\), \(S=V\), self-sponsored disputes, and the hardcoded
\(\mathsf{SHA256}\) case. This made functional testing possible quickly, but it also produced
a monolithic design. The bytecode contained several mutually exclusive branches, so each
deployment paid for code that would never be executed by the selected mode.

The final architecture therefore moved from a monolithic contract to a split architecture.
The selected mode is decided before deployment whenever the mode is known at precontract
time. The contract deployed for one exchange then contains only the logic needed for that
exchange. Dispute-specific choices are delayed until the dispute sponsor roles are known,
because \(S_B=B\) and \(S_V=V\) are not precontract-time facts.

The resulting implementation is not a rewrite of the whole application. It preserves the
existing Rust precontract logic, the Next.js/Electron user interface, the SQLite storage, the
Hardhat test environment, and the ERC-4337-compatible normal path. The main architectural
change is the introduction of narrower Solidity components around the existing execution
flow.

\begin{table}[htbp]
  \centering
  \begin{tabularx}{\linewidth}{>{\raggedright\arraybackslash}p{0.30\linewidth}>{\raggedright\arraybackslash}X}
    \toprule
    Layer & Responsibility \\
    \midrule
    Frontend and backend &
    Select and persist the precontract variant, validate role-dependent choices, and pass
    the selected mode to the deployment layer. \\
    \texttt{SOXFactory} &
    Deploy minimal optimistic clones for specialized precontract modes when reusable
    infrastructure is already available. \\
    Optimistic account implementations &
    Encode the state machine for the honest path and the transition to dispute, with
    variant-specific deposit logic. \\
    Dispute deployers &
    Select the normal, self-sponsored, or hardcoded dispute account when a dispute is
    triggered. \\
    Dispute accounts and verifier libraries &
    Execute Step 7--9, verify commitments, evaluate gates, and implement hardcoded
    \(\mathsf{SHA256}\) reconstruction when selected. \\
    Native Rust CLI tools &
    Generate precontract artifacts and large-file dispute witnesses outside the browser. \\
    \bottomrule
  \end{tabularx}
  \caption{Main layers of the final implementation architecture.}
  \label{tab:architecture-layers}
\end{table}

This split is the key difference between the final implementation and the earlier
monolithic version. The protocol did not become less expressive; rather, the implementation
stopped carrying unrelated modes inside the same deployed bytecode.

\section{Frontend and backend mode selection}

The precontract-time variants are exposed to the application through a small TypeScript
module, \path{src/app/lib/protocol-variants.ts}. It defines four user-facing modes:
\texttt{normal}, \texttt{no\_S\_deposit}, \texttt{S\_equals\_B}, and
\texttt{S\_equals\_V}. The labels shown in the interface are deliberately simpler:
``normal'', ``simSOX (no S-deposit)'', ``S = B'', and ``S = V''.

The vendor selects the mode when creating a precontract in
\path{src/app/components/user/NewContractModal.tsx}. The selected value is sent to
\path{/api/precontracts} either as multipart form data in the web mode or as JSON in the
Electron mode. The backend normalizes the value before inserting it into SQLite. The
\texttt{contracts} table therefore contains a \texttt{precontract\_variant} column, with
\texttt{normal} as the default for backward compatibility.

The same value is later used by the sponsor-facing view. In
\path{src/app/components/sponsor/SponsorContractsListView.tsx}, the interface displays the
mode and enforces the obvious role constraints before deployment:

\begin{itemize}
  \item if the mode is \(S=B\), the deployer must be the buyer;
  \item if the mode is \(S=V\), the deployer must be the vendor;
  \item otherwise, the sponsor can be selected as in the normal flow.
\end{itemize}

These frontend checks are usability checks, not the security boundary. The Solidity factory
and the specialized contracts repeat the relevant role checks on-chain. This duplication is
intentional. It gives a clear error to the user early, while still ensuring that a malformed
API request cannot deploy a contract under an invalid role equality.

The hardcoded \(\mathsf{SHA256}\) circuit variant is not exposed as a normal dropdown option
in the main precontract modal. In the current codebase it is a dedicated deployment and
benchmark path, enabled by passing \texttt{hardcodedSha256Circuit} metadata to the blockchain
deployment helper. This keeps the user-facing mode selection focused on the three
sponsoring simplifications that were requested for the interface: simSOX, \(S=B\), and
\(S=V\).

\section[Factory and optimistic clones]{SOXFactory and clone-based optimistic contracts}

The contract \path{src/hardhat/contracts/SOXFactory.sol} is the main on-chain component of
the clone architecture. It stores four immutable implementation addresses, one for each
optimistic mode: normal, no-deposit, sponsor-is-buyer, and sponsor-is-vendor. For each new
exchange, the factory creates an EIP-1167-style minimal proxy and calls
\texttt{initialize} on the clone.

The factory exposes one creation function per precontract mode:

\begin{itemize}
  \item \texttt{createNormal} for the normal sponsored clone;
  \item \texttt{createNoSDeposit} for simSOX;
  \item \texttt{createSponsorIsBuyer} for \(S=B\);
  \item \texttt{createSponsorIsVendor} for \(S=V\).
\end{itemize}

The functions are not only dispatchers. They also enforce mode-specific constraints. The
no-deposit creation function rejects any initial value, the buyer-sponsor creation function
requires the caller to be the buyer, and the vendor-sponsor creation function requires the
caller to be the vendor. These checks make the selected mode part of the deployed contract's
on-chain history.

In the application deployment helper, the factory is used when clone-based deployment is
enabled, a factory address is present, and the selected mode is a non-hardcoded specialized
variant. The normal interactive path is kept as a direct account deployment because it
preserves the existing ERC-4337-compatible normal contract. The clone factory is
nevertheless available for normal-mode measurements and for comparing direct deployments
with clone-based deployments.

This distinction is important for gas accounting. The factory and implementation contracts
are infrastructure costs: they are deployed once and can then be reused. The marginal cost
of a new clone-based exchange is the cost of creating and initializing the minimal proxy, not
the cost of deploying the full implementation again. This is why later experimental tables
separate one-time infrastructure gas from per-exchange gas.

\section{Specialized optimistic escrow contracts}

The optimistic contract family has two forms: direct contracts and clone contracts. They
share the same high-level state machine, but they are used in different integration contexts.

The direct variants extend a common direct base. They are full contract deployments with
constructor arguments. They are useful as fallbacks and are used by the hardcoded
\(\mathsf{SHA256}\) path. The direct normal contract is also the path that keeps the
ERC-4337-style account functionality inherited from the original application.

The clone variants extend \path{OptimisticSOXAccountCloneBase.sol}. They are initialized
after creation by the factory. The clone base removes the ERC-4337 account logic and keeps
only the optimistic escrow behavior needed by the specialized variants. This was a
deliberate simplification: the variants are measured and demonstrated mainly with direct
transactions, while the normal path remains available when the existing account-abstraction
flow is needed.

The four clone implementations encode the deposit differences described in
Chapter~\ref{chap:variants}. The normal clone requires the sponsor fee and waits for buyer
payment. The no-deposit clone requires zero value and makes the buyer payment equal to the
agreed price only. The \(S=B\) clone requires the buyer to be the deployer, records the buyer
payment during initialization, and starts directly in the key-waiting state. The \(S=V\)
clone requires the vendor to be the deployer, records the sponsor deposit, and still waits
for buyer payment.

Both direct and clone bases contain the same dispute-sponsor entry points needed after the
honest path fails. In particular, they include:

\begin{itemize}
  \item a self-sponsored buyer dispute deposit path;
  \item an authorization-based external buyer dispute sponsor path;
  \item a self-sponsored vendor dispute deposit path;
  \item a runtime choice between normal and self-sponsored dispute deployment.
\end{itemize}

The authorization-based buyer path is especially important. It prevents an external
\(S_B\) from unilaterally claiming that the buyer is unhappy. The signed message includes
the chain id, the optimistic contract address, the buyer, the proposed buyer-side dispute
sponsor, and the commitment. This binds the authorization to one exchange and one proposed
sponsor.

\section{Specialized dispute game contracts}

The dispute architecture is split into deployer libraries and dispute account contracts.
The deployer libraries are small selectors that instantiate the appropriate dispute account:
normal, self-sponsored, or hardcoded \(\mathsf{SHA256}\). The optimistic account decides
which deployer to use when the vendor-side dispute sponsor deposit is made.

The normal dispute account, \path{DisputeSOXAccountNormal.sol}, keeps the generic dispute
behavior. It implements the challenge-response loop, Step 8 verification, and the normal
Step 9 behavior. This path is used when the dispute sponsors are external or when the
self-sponsored specialization is not applicable.

The self-sponsored dispute account,
\path{DisputeSOXAccountSelfSponsored.sol}, inherits from the normal dispute account and
overrides only the Step 9 handling. When both \(S_B=B\) and \(S_V=V\), it does not relaunch
the challenge loop for a new sponsor configuration. Instead, it moves directly toward the
terminal outcome determined by the losing party. This keeps the implementation small and
makes the optimization auditable: the rest of the dispute logic is inherited unchanged, and
the only semantic change is the Step 9 transition.

This inheritance structure was chosen to reduce implementation risk. The project did not
rewrite the complete dispute verifier for the self-sponsored case. It specialized the part
that is actually different. This also makes the gas measurements easier to interpret,
because a saving in the self-sponsored case can be attributed to removing unnecessary
iterations rather than to a hidden rewrite of the verifier.

\section{Hardcoded SHA-256 dispute contracts}

The hardcoded \(\mathsf{SHA256}\) path is implemented separately from the clone factory. Its
optimistic contract is \path{OptimisticSOXAccountHardcodedSHA256.sol}. This contract stores
three immutable pieces of metadata: the description hash, the plaintext length, and the
ciphertext IV. At deployment time, it verifies that the number of ciphertext blocks and the
number of circuit gates match the hardcoded \(\mathsf{SHA256}\) layout.

When a dispute is triggered from this contract, it uses the hardcoded SHA-256 dispute
deployer, which instantiates the corresponding hardcoded dispute account. The dispute
account overrides the generic Step 8 entry points so that the vendor no longer provides the
circuit gate bytes as the source of truth and no longer provides the \(\pi_1\) Merkle proof
for \(h_{\mathit{circuit}}\). In the strict hardcoded tests, the corresponding inputs are
empty.

The reconstruction logic is isolated in
\path{HardcodedSha256CircuitLib.sol}. Given the gate number, the number of blocks, the
description hash, the plaintext length, and the IV, the library computes the expected
64-byte encoding of the selected gate and checks its hash. It covers the AES-CTR input
gates, the SHA-256 padding gates, the SHA-256 compression chain, the description constant,
and the final comparison gate.

The hardcoded contract does not remove all witness data from Step 8. The smart contract
does not store the whole ciphertext nor every intermediate value. The vendor still provides
values needed to evaluate the selected gate, but those values are verified against the
appropriate commitments and accumulator proofs before the gate result is accepted. The
architectural change is therefore specific and limited: \(g_i\) is reconstructed by the
contract, while \(v_i\)-style witness values remain externally supplied and authenticated.

This design also explains why hardcoded \(\mathsf{SHA256}\) is implemented as a dedicated
path instead of a flag inside every dispute contract. The hardcoded path changes the Step 8
verifier. Carrying both generic and hardcoded verification in the same contract would put
the project back into the monolithic pattern that the final architecture was meant to avoid.

\section{Native CLI path for large files}

The original application relied heavily on WebAssembly and browser-adjacent execution. That
is convenient for small demonstrations, but it becomes fragile for very large files. The
project therefore uses native Rust command-line tools for large-file experiments.

The main precontract tool is \path{src/wasm/src/bin/precontract_cli.rs}. It reads an input
file, chooses or receives a 16-byte AES key, computes the V2 precontract values, writes the
ciphertext and circuit bytes to disk, and returns a JSON object containing:
\texttt{description\_hex}, \texttt{h\_ct\_hex}, \texttt{h\_circuit\_hex},
\texttt{commitment\_c\_hex}, \texttt{commitment\_o\_hex}, \texttt{num\_blocks},
\texttt{num\_gates}, output paths, and the key.

The web API uses this CLI in multipart mode. Instead of loading the upload through a
browser-only computation path, \path{/api/precontracts} streams the uploaded file to a
temporary file and calls \texttt{precontract\_cli}. In Electron mode, the application can run
the same native precomputation locally and send the resulting JSON to the backend. In both
cases, the backend stores the precontract metadata and moves the ciphertext artifact to the
application upload directory.

The large-file dispute benchmark uses a second native tool,
\path{src/wasm/src/bin/hardcoded_sha256_dispute_cli.rs}. This tool is intentionally narrower
than a full generic dispute prover. It validates the number of blocks and gates for the
hardcoded \(\mathsf{SHA256}\) layout, computes the sequence of \(hpre_i\) values for the
left-branch benchmark path, and prepares the final left-boundary Step 8 witness. This is the
tool used to make the 1 GiB hardcoded dispute measurements practical.

The distinction between the two native tools matters. The precontract CLI is part of the
interactive creation path. The hardcoded dispute CLI is a benchmark and measurement tool for
a selected dispute scenario. It does not mean that the full browser UI now supports every
possible 1 GiB generic dispute interaction end-to-end.

\section{Database and artifact management}

The application stores precontract and exchange metadata in SQLite through
\path{src/app/lib/sqlite.ts}. The schema is defined in \path{src/app/db/init.sql} and
contains two main tables: \texttt{contracts} and \texttt{disputes}. The project added the
\texttt{precontract\_variant} column to \texttt{contracts}, with default value
\texttt{normal}. The runtime initialization code also attempts to add this column to older
local databases, which keeps existing development databases usable.

The \texttt{contracts} table stores the buyer and vendor addresses, price, tips, timeout,
number of blocks, number of gates, commitment, opening value, protocol version, algorithm
suite, acceptance status, sponsor address, deployed optimistic contract address, and session
key information. The \texttt{disputes} table stores the dispute sponsors and dispute
contract address for a contract id.

Encrypted files are stored outside the database. In web mode, the backend writes the upload
to a temporary file, invokes the native precontract CLI, and then copies the produced
ciphertext to the uploads directory under a deterministic name based on the contract id. In
Electron mode, the precomputed ciphertext path is sent by the desktop shell and uploaded
through the Electron bridge. This design avoids storing large ciphertexts directly in
SQLite.

The backend also returns the key, \(h_{\mathit{circuit}}\), and \(h_{\mathit{ct}}\) to the
frontend after precontract creation. The current application stores some of these values in
browser local storage for the demonstration workflow. This is acceptable for the local proof
of concept, but it is not a production key-management design. The limitation is documented
explicitly later in the report rather than hidden as an implementation detail.

\section{Integration with the existing application}

A central constraint of the project was to improve the implementation without breaking the
existing proof-of-concept application. The final architecture therefore integrates the new
contracts through the existing application lifecycle:

\begin{enumerate}
  \item the vendor creates a precontract and selects a transaction mode;
  \item the backend computes or receives the Rust precontract output and stores metadata;
  \item the buyer accepts the precontract and downloads the ciphertext;
  \item the sponsor view deploys the optimistic contract using the selected mode;
  \item the buyer pays if the selected mode still requires a separate payment;
  \item the vendor releases the key and the exchange completes optimistically, or the parties
  move to the dispute path.
\end{enumerate}

The deployment helper is the bridge between the application and the Solidity architecture.
It reads deployed infrastructure addresses from the generated deployment JSON file, chooses
the factory path when applicable, or falls back to direct deployment when the selected mode
requires it. It also checks for unlinked library placeholders before deploying a direct
contract, which prevents a class of confusing runtime failures where bytecode artifacts were
not regenerated after a contract change.

The normal path remains compatible with the original ERC-4337 and EntryPoint demonstration. The
specialized clone variants do not claim full account-abstraction support; their
\texttt{supportsERC4337} function returns false through the clone base. This separation is
intentional. It lets the project measure the gas benefit of narrower contracts without
pretending that every optimized variant has the same transaction-sponsoring integration as
the inherited normal account.

The final integrated architecture is therefore hybrid but controlled. The user-facing
application can demonstrate normal sponsored exchanges and the main precontract variants.
The benchmark suite can additionally exercise clone-based normal measurements, hardcoded
\(\mathsf{SHA256}\) disputes, and large-file native paths. The report keeps these scopes
separate in the experimental chapters so that an application demo, a gas benchmark, and a
large-file native benchmark are not accidentally treated as the same result.
